% English translation of the front matter of SGA 1 (abstract, preface,
% Introduction, Avertissement), from lines 1--545 of the corrected SMF
% source smf_doc-math_3_01.tex.

\chapter*{Abstract}

This volume is a newly typeset and annotated edition of the book
``Rev\^etements \'Etales et Groupe Fondamental'', Lecture Notes in
Mathematics, 224, Springer-Verlag, Berlin-Heidelberg-New York, 1971, by
Alexander Grothendieck et al.

The text presents the foundations of a theory of the fundamental group
in Algebraic Geometry, from the ``Kroneckerian'' point of view, which
makes it possible to treat on the same footing the case of an algebraic
variety in the usual sense and that of the ring of integers of a number
field, for example.

\medskip
\noindent\emph{Mathematics Subject Classification:} 14-02, 14A15, 14B25,
14D15, 14E20, 14F35, 14H30.

\smallskip
\noindent\emph{Keywords:} \'Etale morphism, smooth morphism, flat
morphism, scheme, fundamental group, covering, theory of descent,
specialization.

\chapter*{Preface}

The present text is a newly typeset and annotated edition of the book
``Rev\^etements \'Etales et Groupe Fondamental'', Lecture Notes in
Mathematics, 224, Springer-Verlag, Berlin-Heidelberg-New York, 1971, by
Alexander Grothendieck et al.

The typesetting in \LaTeX 2e was carried out by volunteers participating
in a project directed by Bas Edixhoven, about which more details can be
found at the address \url{http://www.math.leidenuniv.nl/~edix/}. The
page layout was completed by the \hbox{Soci\'et\'e} math\'ematique de
France.

This is a \emph{slightly corrected version} of the original text. It is
published by the \hbox{Soci\'et\'e} math\'ematique de France
(\url{http://smf.emath.fr/}) in the series ``Documents
\hbox{Math\'ematiques}''. Some updates have been made by Michel Raynaud.
These are delimited by brackets \lcrochetbf\,\rcrochetbf\ and indicated
by the symbol (MR)
% source: page references replaced by statement numbers
(Remarks X~\Ref{X.2.14}, XI~\Ref{XI.1.4}, XII~\Ref{XII.5.6},
XIII~\Ref{XIII.2.13}, and the footnote to III~6.6).
In order to unify notation, the residue field of a point $x$ is denoted
$\kres(x)$, and that of a local ring $A$ is denoted $\kres(A)$.\par
There also exists an electronic version meant to reproduce the original
text.

The two versions have a single common source file, which is to be found
on the arXiv.org e-Print archive server at \url{http://arxiv.org/}. The
differences between the two versions are documented in the source file.

The old page numbering is incorporated in the margin (the number~$n$
marking the beginning of page~$n$).\footnote{Translator's note: in this
English edition the original page numbers are kept only as comments in
the \TeX{} source, and the page references of the SMF edition in the
preceding paragraphs have been replaced by statement numbers.}

\chapter*{Introduction}
\setcounter{section}{0}\renewcommand{\thesubsection}{\arabic{subsection}}
\label{I.0}

In the first part of this introduction we give some details on the
contents of the present volume; in the second, on the whole of the
``\emph{S\'eminaire de G\'eom\'etrie Alg\'ebrique du Bois-Marie}'', of
which the present volume constitutes the first volume.

\subsection{}
\label{I.introduction.1}
The present volume presents the foundations of a theory of the
fundamental group in Algebraic Geometry, from the ``Kroneckerian'' point
of view, which makes it possible to treat on the same footing the case of
an algebraic variety in the usual sense and that of the ring of integers
of a number field, for example. This point of view can be expressed in a
satisfactory way only in the language of schemes, and we shall use this
language freely, as well as the main results expounded in the first
three chapters of the \emph{\'El\'ements de G\'eom\'etrie Alg\'ebrique}
of J\ptbl \textsc{Dieudonn\'e} and A\ptbl \textsc{Grothendieck} (cited
EGA in what follows). The study of the present volume of the
``\emph{S\'eminaire de G\'eom\'etrie Alg\'ebrique du Bois-Marie}''
requires no other knowledge of Algebraic Geometry, and can therefore
serve as an introduction to the present-day techniques of Algebraic
Geometry for a reader wishing to become familiar with these techniques.

Expos\'es \Ref{I} to \Ref{XI} of this book are a verbatim reproduction,
practically unchanged, of the mimeographed notes of the oral Seminar,
which were distributed by the \emph{Institut des Hautes \'Etudes
Scientifiques}\footnote{As were the notes of the seminars following this
one. This mode of distribution having proved impracticable and
insufficient in the long run, all the ``\emph{S\'eminaire de
G\'eom\'etrie Alg\'ebrique du Bois-Marie}'' will henceforth appear in
book form, like the present volume.}. We have confined ourselves to
adding a few footnotes to the original text, correcting some misprints,
and making a terminological adjustment, the term ``simple morphism'' in
particular having been replaced in the meantime by ``smooth morphism'',
which does not lend itself to the same confusions.

Expos\'es \Ref{I} to \Ref{IV} present the local notions of \emph{\'etale}
morphism and \emph{smooth} morphism; they make hardly any use of the
language of schemes, expounded in Chapter~I of the
\emph{\'El\'ements}\footnote{A more complete study is now available in
the \emph{\'El\'ements}, Chap.~IV, \oldS\S 17 and 18.}.
Expos\'e~\Ref{V} presents the axiomatic description of the fundamental
group of a scheme, useful even in the classical case where this scheme
reduces to the spectrum of a field, where one finds a very convenient
reformulation of the usual Galois theory. Expos\'es \Ref{VI} and
\Ref{VIII} present the \emph{theory of descent}, which has taken on a
growing importance in Algebraic Geometry in recent years, and which could
render analogous services in Analytic Geometry and in Topology. It
should be pointed out that Expos\'e VII had not been written up, and its
substance is to be found incorporated in a work of J\ptbl Giraud
(M\'ethode de la Descente, Bull.\ Soc.\ Math.\ France, M\'emoire~2,
1964, viii + 150~p.). In Expos\'e~\Ref{IX} one studies more specifically
the descent of \'etale morphisms, obtaining a systematic approach to
theorems of \textsc{Van Kampen} type for the fundamental group, which
appear here as simple translations of descent theorems. It is
essentially a procedure for computing the fundamental group of a
connected scheme~$X$ endowed with a surjective and proper morphism, say
$X'\to X$, in terms of the fundamental groups of the connected
components of~$X'$ and of the fibered products~$X'\times_X X'$,
$X'\times_X X'\times_X X'$, and of the homomorphisms induced between
these groups by the canonical simplicial morphisms between the preceding
schemes. Expos\'e~\Ref{X} gives the theory of \emph{specialization of the
fundamental group} for a proper and smooth morphism, the most striking
result of which consists in the determination (to within very little) of the
fundamental group of a smooth algebraic curve in characteristic $p>0$,
thanks to the result known by transcendental means in characteristic
zero. Expos\'e~\Ref{XI} gives some examples and complements, making
explicit in particular, in cohomological form, the theory of
\emph{\textsc{Kummer}} coverings, and that of
\emph{\textsc{Artin}-\textsc{Schreier}}. For further comments
on the text, see the \emph{Foreword} to the mimeographed version, which
follows the present Introduction.

Since the writing of the present Seminar in 1961, there have been
developed, in collaboration by M\ptbl \textsc{Artin} and myself, the
language of \emph{\'etale topology} and a corresponding cohomological
theory, expounded in the part SGA~4 ``\'Etale cohomology of schemes'' of
the \emph{S\'eminaire de G\'eom\'etrie Alg\'ebrique}, to appear in the
same series as the present volume. This language, and the results it
already has at its disposal, provide a particularly flexible tool for the
study of the fundamental group, making it possible to understand better,
and to go beyond, some of the results expounded here. There would
therefore be reason to take up the theory of the fundamental group
entirely anew from this point of view (all the key results in fact
already appearing in \loccit). This is what was planned for the chapter
of the \emph{\'El\'ements} devoted to the fundamental group, which was
also to contain several other developments that could not find their
place here (relying on the technique of resolution of singularities):
computation of the ``local fundamental group'' of a complete local ring
in terms of a suitable resolution of singularities of this ring, local
and global K\"unneth formulas for the fundamental group without a
properness hypothesis (\cf Exp\ptbl \Ref{XIII}), the results of
M\ptbl \textsc{Artin} on the comparison of the local fundamental groups
of an excellent henselian local ring and of its completion
(SGA~4~XIX). Let us also point out the necessity of developing a theory
of the fundamental group of a topos, which will encompass at the same
time the ordinary topological theory, its semi-simplicial version, the
``profinite'' variant developed in Expos\'e~\Ref{V} of the present volume,
and the somewhat more general pro-discrete variant of SGA~3~X~7 (adapted
to the case of non-normal and non-unibranch schemes). Pending an overall
recasting of the theory from this point of view, Expos\'e \Ref{XIII} by
Mme \textsc{Raynaud}, which uses the language and the results of SGA~4,
is intended to show the advantage that can be drawn from them in some
typical questions, in particular by generalizing certain results of
Expos\'e~\Ref{X} to non-proper relative schemes. In it one gives in
particular the structure of the ``prime-to-$p$'' fundamental group of a
non-complete algebraic curve in arbitrary characteristic (which I had
announced in 1959, but of which no proof had been published up to now).

In spite of these numerous gaps and imperfections (others will say:
because of these gaps and imperfections), I think that the present
volume may be useful to the reader who wishes to become familiar with
the theory of the fundamental group, as well as a reference work,
pending the writing and publication of a text escaping the criticisms I
have just enumerated.

\subsection{}
\label{I.introduction.2}
The present volume constitutes volume 1 of the ``\emph{S\'eminaire de
G\'eom\'etrie Alg\'ebrique du Bois-Marie}'', the following volumes of
which are scheduled to appear in the same series as this one. The aim
that the \emph{Seminar} sets itself, in parallel with the treatise
``\emph{\'El\'ements de G\'eom\'etrie Alg\'ebrique}'' of
J\ptbl \textsc{Dieudonn\'e} and A\ptbl \textsc{Grothendieck}, is to lay
the foundations of Algebraic Geometry, following the points of view in
this latter work. The standard reference for all the volumes of the
\emph{Seminar} is constituted by Chapters~I, II, III of the
``\emph{\'El\'ements de G\'eom\'etrie Alg\'ebrique}'' (cited EGA~I, II,
III), and the reader is assumed to be in possession of the background in
commutative algebra and homological algebra that these chapters
presuppose\footnote{See the Introduction to EGA~I for details on this
point.}. Moreover, in each volume of the \emph{Seminar} reference will
be made freely, to the extent needed, to earlier volumes of the same
\emph{Seminar}, or to other chapters of the ``\'El\'ements'' that have
been published or are about to appear.

Each \emph{part} of the \emph{Seminar} is centered on a principal
subject, indicated in the title of the corresponding volume or volumes;
the oral seminar generally extends over one academic year, sometimes
more. The expos\'es within each part of the \emph{Seminar} are generally
in close logical dependence on one another; on the other hand, the
different parts of the \emph{Seminar} are to a large extent logically
independent of one another. Thus the part ``Group Schemes'' is almost
entirely independent of the two parts of the \emph{Seminar} that precede
it chronologically; on the other hand, it makes frequent appeal to the
results of EGA~IV. Here is the list of the parts of the \emph{Seminar}
that are to appear shortly (cited SGA~1 to SGA~7 in what follows):
\begin{enumerateb}
\item[SGA 1.] \'Etale coverings and the fundamental group, 1960 and 1961.
\item[SGA 2.] Local cohomology of coherent sheaves and local and global
Lefschetz theorems, 1961/62.
\item[SGA 3.] Group schemes, 1963 and 1964 (3~volumes), in~coll. with
M\ptbl \textsc{Demazure}.
\item[SGA 4.] Theory of topoi and \'etale cohomology of schemes, 1963/64
(3~volumes), (in~coll. with M\ptbl \textsc{Artin} and
J.L\ptbl \textsc{Verdier}).
\item[SGA 5.] $l$-adic cohomology and $L$-functions, 1964 and 1965
(2~volumes).
\item[SGA 6.] Theory of intersections and the Riemann-Roch theorem,
1966/67 (2~volumes) (in~coll. with P\ptbl \textsc{Berthelot} and
L\ptbl \textsc{Illusie}).
\item[SGA 7.] Local monodromy groups in algebraic geometry.
\end{enumerateb}

Three of these partial \emph{Seminars} were directed in
\emph{collaboration} with other mathematicians, who will appear as
co-authors on the cover of the corresponding volumes. As for the other
active participants of the \emph{Seminar}, whose role (as much from the
point of view of writing up as from that of the work of mathematical
refinement) has been growing from year to year, the name of each
participant appears at the head of the expos\'es for which he is
responsible as lecturer or as writer, and the list of those who appear
in a given volume is indicated on the flyleaf of the said volume.

It is appropriate to give some details on the relations between the
\emph{Seminar} and the \emph{\'El\'ements}. The latter were intended in
principle to give a comprehensive exposition of the notions and
techniques judged most fundamental in Algebraic Geometry, as these
notions and techniques themselves emerge, through the natural play of
requirements of logical and aesthetic coherence. From this perspective,
it was natural to regard the \emph{Seminar} as a preliminary version of
the \emph{\'El\'ements}, destined to be absorbed almost entirely, sooner
or later, into the latter. This process had already begun to a certain
extent some years ago, since Expos\'es I to IV of the present volume SGA~1
are entirely absorbed by EGA~IV, and Expos\'es VI to VIII were to be
absorbed within a few years in EGA~VI. However, as the work of
construction undertaken in the \emph{\'El\'ements} and in the
\emph{Seminar} develops, and as its overall proportions become clearer,
the initial principle (according to which the \emph{Seminar} would
constitute only a preliminary and provisional version) appears less and
less realistic because of (among other things) the limits imposed by
provident nature on the length of human life. Taking into account the
care generally taken in writing up the various parts of the
\emph{Seminar}, there will doubtless be reason to take up such a part
again in the \emph{\'El\'ements} (or in treatises that might take over
from them) only when progress subsequent to its writing makes it
possible to bring very substantial improvements to it, at the price of
fairly deep modifications. This is already the case for the present
seminar SGA~1, as was said above, and also for SGA~2 (thanks to the
recent results of Mme~\textsc{Raynaud}). On the other hand, nothing
indicates at present that this will be so in the near future for any of
the parts SGA~3 to SGA~7 cited above.

References within the ``S\'eminaire de G\'eom\'etrie Alg\'ebrique du
Bois Marie'' are given as follows. An \emph{internal reference} to one
of the parts SGA~1 to SGA~7 of the Seminar is given in the style
III~9.7, where the numeral III denotes the number of the expos\'e, which
appears at the top of each page of the expos\'e in question, and~9.7 the
number of the statement (or of the definition, remark, etc.) within the
expos\'e. Where appropriate, longer decimal numbers may be used, for
example 9.7.1, 9.7.2 to denote, for example, the various steps in the
proof of a proposition~9.7. The reference III~9 denotes paragraph 9 of
Expos\'e~III. The number of the expos\'e is omitted for references internal
to an expos\'e. For a \emph{reference to another} of the \emph{parts} of
the \emph{Seminar}, the same abbreviations are used, but preceded by mention
of the part of the SGA in question, SGA~1~III~9.7. Likewise, the
reference EGA~IV~11.5.7 means: \'El\'ements de G\'eom\'etrie
Alg\'ebrique, Chap\ptbl IV, statement (or definition etc.)~11.5.7; here
the first arabic numeral again denotes the number of the paragraph.
Apart from these conventions, in force throughout the SGA, the
bibliography relating to an expos\'e will generally be gathered at the end
of it, and will be referred to within the expos\'e by numbers in square
brackets such as [\textbf{3}], following the usual practice.

Finally, for the convenience of the reader, whenever it seems necessary,
we shall append at the end of the volumes of the SGA an index of
notation and a terminological index containing, where appropriate, an
English translation of the French terms used.

I wish to add to this introduction an extra-mathematical comment. In the
month of November 1969 I learned that the Institut des Hautes \'Etudes
Scientifiques, of which I have been a professor essentially since its
foundation, had for three years been receiving subsidies from the
Ministry of the Armed Forces. Already as a beginning researcher I found
extremely regrettable the lack of scruples that most scientists show in
agreeing to collaborate in one form or another with the military
apparatuses. My motivations at that time were essentially of a moral
nature, and therefore hardly likely to be taken seriously. Today they
acquire a new force and a new dimension, in view of the danger of
destruction of the human species with which we are threatened by the
proliferation of the military apparatuses and of the means of mass
destruction that these apparatuses have at their disposal. I have
explained myself elsewhere in more detail on these questions, which are
much more important than the advancement of any science whatsoever
(mathematics included); one may for example consult on this subject the
article by G\ptbl Edwards in \no 1 of the journal Survivre (August 1970),
summarizing a more detailed exposition of these questions that I had
given elsewhere. Thus I found myself working for three years in an
institution while it was taking part, without my knowledge, in a mode of
financing that I consider immoral and dangerous\footnote{It goes without
saying that the opinion I have just expressed is my own
responsibility alone, and not that of the publishing house Springer,
which publishes the present volume.}. Being at present alone in holding this
opinion among my colleagues at the IHES, which has doomed to failure my
efforts to obtain the removal of the military subsidies from the budget
of the IHES, I have taken the decision that was called for, and am
leaving the IHES on September 30, 1970, and am likewise suspending all
scientific collaboration with this institution for as long as it
continues to accept such subsidies. I have asked Mr.~Motchane, director
of the IHES, that from October 1, 1970 on the IHES refrain from
distributing mathematical texts of which I am the author, or which form
part of the S\'eminaire de G\'eom\'etrie Alg\'ebrique du Bois Marie. As
was said above, the distribution of this seminar will be taken over by
the firm Julius Springer, in the series of the Lecture Notes. I am happy
to thank here the Springer firm and Mr.~K\ptbl Peters for the efficient
and courteous help they have given me in making this publication
possible, by taking charge in particular of the typing, for photo-offset,
of the new expos\'es added to the old seminars, and of the missing expos\'es
of the incomplete seminars.

I also thank Mr.~J.P\ptbl Delale, who took on the thankless task of
compiling the index of notation and the terminological index.

\bigskip\hfill Massy, August 1970.

\renewcommand{\thesubsection}{\thesection.\arabic{subsection}}

\chapter*{Foreword}
\label{I.avertissement}

Each of the written-up expos\'es gives the substance of several
consecutive oral expos\'es. It did not seem useful to specify their dates.

Expos\'e~VII, to which reference is made at various points in
Expos\'e~\Ref{VIII}, has not been written up by the lecturer, who in the
oral lectures had confined himself to sketching the language of descent
in general categories, taking a strictly utilitarian point of view and
without going into the logical difficulties raised by this language. It
became apparent that a correct exposition of this language would go
beyond the limits of the present notes, if only by its length. For a
formal exposition of the theory of descent, I refer to an article in
preparation by Jean \textsc{Giraud}. Until it
appears\footnote{\label{footnotegiraud}Now published:
J\ptbl \textsc{Giraud}, \emph{M\'ethodes de la descente}, M\'emoire
\no 2 of the Soci\'et\'e math\'ematiques de France, 1964.}, I think that
an attentive reader will have no trouble in supplying by his own means
the phantom references of Expos\'e~VIII.

Other oral expos\'es, coming after Expos\'e~\Ref{XI}, and to which allusion
is made at certain places in the text, have not been written up either,
and were intended to form the substance of an Expos\'e~\Ref{XII} and of an
Expos\'e~\Ref{XIII}. The first ones among these oral expos\'es took up again, in the
framework of schemes and of analytic spaces with nilpotent elements (as
they are introduced in the S\'eminaire Cartan 1960/61), the
construction of the analytic space associated with a prescheme locally
of finite type over a complete valued field~$k$, the theorems of GAGA
type in the case where $k$ is the field of complex numbers, and the
application to the comparison of the fundamental group defined by
transcendental means and the fundamental group studied in these notes
(compare A\ptbl Grothendieck, Fondements de la G\'eom\'etrie
Alg\'ebrique, S\'eminaire Bourbaki \no 190, page~10, December 1959). The
last oral expos\'es sketched the generalization of the methods developed
in the text to the study of coverings admitting tame ramification, and
of the structure of the fundamental group of a complete curve with a
finite number of points removed (compare \loccit, \no 182, page 27,
theorem~14). These expos\'es introduce no essentially new idea, which is
why it did not seem indispensable to give a formal write-up of them
before the appearance of the corresponding chapters of the \'El\'ements
de G\'eom\'etrie Alg\'ebrique\footnote{They are included in the present
volume in Exp\ptbl \Ref{XII} by Mme~Raynaud, with a proof different from
the original proof expounded in the oral Seminar (\cf introduction).}.

On the other hand, the theorems of Lefschetz type for the fundamental
group and the Picard group, from the local as well as from the global
point of view, have been the subject of a separate Seminar in 1962,
which has been completely written up and is available to
users\footnote{\emph{\'Etale cohomology of coherent sheaves and local and
global Lefschetz theorems} (SGA~2), published by North Holland Pub.\
Co.}. Let us point out that the results developed both in the present
Seminar and in that of 1962 will be used in an essential way in the
appearance of several key results in the \'etale cohomology of preschemes,
which will be the subject of a Seminar (conducted by M\ptbl Artin and
myself) in 1963/64, at present in preparation\footnote{\emph{\'Etale
cohomology of schemes} (cited SGA~4), to appear in this same series.}.

Expos\'es \Ref{I} to~\Ref{IV}, of an essentially local and very elementary
nature, will be entirely absorbed by Chapter IV of the
\emph{\'El\'ements de G\'eom\'etrie Alg\'ebrique}, the first part of
which is in press and which will doubtless be published around the end
of '64. They may nevertheless be useful to a reader who would like to
become acquainted with the essential properties of smooth, \'etale or flat
morphisms before entering into the arcana of a systematic treatise. As
for the other expos\'es, they will be absorbed into Chapter
VIII\footnote{In fact, as a result of a modification of the plan
initially envisaged for the \emph{\'El\'ements}, the study of the
fundamental group is postponed there to a chapter later than the one
just indicated. Compare the introduction that precedes the present
``Foreword''.} of the ``\'El\'ements'', whose publication can hardly be
envisaged for several years.

\bigskip\hfill Bures, June 1963.
